\documentclass[10pt,a4paper]{amsart}
\usepackage[margin=19mm]{geometry}
\usepackage{amsmath,amssymb,mathtools}
\usepackage[scr=rsfs,cal=euler]{mathalfa}
\usepackage{xfrac}
\usepackage[unicode,hidelinks,bookmarks=false]{hyperref}
\newcommand{\ie}{i.e.\ }
\newcommand{\cf}{cf.\ }
\newcommand{\resp}{respectively\ }
\newcommand{\Subsection}{\subsection}
\providecommand{\nobreakdash}{\nobreak\hbox{-}}
\setlength{\emergencystretch}{2em}
\multlinegap0pt
\usepackage{tikz}
\usetikzlibrary{cd}
\usepackage[shortlabels]{enumitem}
\usepackage{mathtools}
\usepackage{amssymb}
\usepackage{xcolor}
\newtheorem{theorem}{Theorem}
\newtheorem{claim}{Claim}
\newcommand{\CP}[1]{{#1}^{\text{\tiny$\vee$}}}			%Small copresheaves on #1
\newcommand{\ES}{\textnormal{ES}}
\newcommand{\IFF}{\Leftrightarrow}						% Iff
\newcommand{\IMP}{\Rightarrow}						% Left-to-right implication
\newcommand{\Set}{{\sf Set}}							%Sets
\newcommand{\cat}{{\sf C}}								%A generic category 
\renewcommand{\epsilon}{\varepsilon}
\newcommand{\op}{{\mathrm{op}}}						% Opposite category
\renewcommand{\phi}{\varphi}
\newcommand{\yo}{\mathcal{Y}}
\title{On implicit operations and balanced categories}
\author{Luca Reggio}
\date{Source: arXiv:2609.15626v1 (CC BY 4.0)}
\begin{document}
\maketitle

\noindent\emph{Source and licence.} On implicit operations and balanced categories. Version arXiv:2609.15626v1 (CC BY 4.0), distributed by arXiv under the Creative Commons Attribution 4.0 International licence, http://creativecommons.org/licenses/by/4.0/, as recorded in the arXiv licence metadata for this exact version and shown on the abstract page https://arxiv.org/abs/2609.15626v1. Full licence notice: \texttt{licenses/arxiv-2609.15626-CC-BY-4.0.txt}.\par\smallskip

\noindent\textit{Editorial scope.} Counted unit: the article's theorem t:ES-vs-implicit-operations (source0.tex lines 344--351). The article's complete original proof of it is delivered with it (source0.tex lines 353--395, one \texttt{proof} environment of 43 lines). No mathematics is written, repaired or re-derived: the statement and the proof are the article's own lines, in the article's own notation, and no retained line is substituted or re-flowed.  No further source-local context is needed: every cross-reference of every retained line resolves inside the two delivered spans.  Nothing was omitted from any retained span, so the package carries no omission record.  No retained line contains a citation, so the wrapper carries no bibliography and no bibliography entry.  The retained lines contain the article's own diagram code, which the wrapper typesets with the standard tikz packages; no image file is delivered and the package declares no asset dependency.  Editorial formatting only, in addition: an AMS \texttt{amsart} preamble that restates the article's own declarations and macros used by the retained lines, namely the environments \texttt{theorem}, \texttt{claim} on separate counters, together with the standard packages those lines need.  The retained environments print in this document as Theorem 1, Claim 1 in order of appearance, whereas the article prints the same items with the numbering of its own full document.  The complete native source of the article as distributed in the arXiv e-print for this exact version is bundled unchanged as \texttt{sources/210433/source0.tex}.
\begin{theorem}\label{t:ES-vs-implicit-operations}
Let $\cat$ be a complete category with an extremal projective generator~$G$, and let $V\coloneqq \cat(G,-)\colon\cat\to\Set$. The following statements are equivalent:
%
\begin{enumerate}[label=(\arabic*),series=implicitops]
\item\label{i:Cat-satisfies-ES} $\cat$ satisfies \ES~(equivalently, $\cat$ is balanced);
\item\label{i:representable-admits-extension} every representable implicit partial operation $V^{I} \rightsquigarrow V$ can be extended to a total one.
\end{enumerate}
\end{theorem}

\begin{proof}
Before starting, let us mention in passing that the implication \ref{i:Cat-satisfies-ES} $\IMP$ \ref{i:representable-admits-extension} holds even when $G$ is not an extremal generator, and the implication \ref{i:representable-admits-extension} $\IMP$ \ref{i:Cat-satisfies-ES} holds even when $G$ is not extremal projective. 

We show that both items~\ref{i:Cat-satisfies-ES} and~\ref{i:representable-admits-extension} are equivalent to the following: 
\begin{enumerate}[label=(\arabic*),resume=implicitops]
\item\label{i:ES-at-copowers-of-G} every epimorphism in $\cat$ whose domain is a small copower of $G$ is an extremal epimorphism.
\end{enumerate}

Clearly, \ref{i:Cat-satisfies-ES} $\IMP$ \ref{i:ES-at-copowers-of-G}. Conversely, suppose that $\delta\colon A\to B$ is an epimorphism in $\cat$. Let~$I$ be the set $\cat(G,A)$, and consider the canonical arrow $\epsilon\colon \coprod_{I}G \to A$. The latter is an (extremal) epimorphism because $G$ is an extremal generator. The composite $\delta\circ \epsilon \colon \coprod_{I}G \to B$ is an epimorphism whose domain is a small copower of $G$ and thus, by assumption, it is an extremal epimorphism. It follows easily that $\delta$ is an extremal epimorphism too.

It remains to show that \ref{i:representable-admits-extension} $\IFF$ \ref{i:ES-at-copowers-of-G}. We start by proving the following fact:
%
\begin{claim}\label{claim:ext-epi-iff-total}
Let $I$ be a set and let $\delta\colon \coprod_{I}G \to A$ be an epimorphism in $\cat$. Then~$\delta$ is an extremal epimorphism if, and only if, every representable implicit partial operation $V^{I} \rightsquigarrow V$ with domain $\yo(\delta)\colon \yo(A)\rightarrowtail V^{I}$ can be extended to a total one.
\end{claim}
%
\begin{proof}[Proof of Claim]
Assume that $\delta$ is an extremal epimorphism and let $f\colon \yo(A)\to V$ be a representable implicit partial operation with domain $\yo(\delta)\colon \yo(A)\rightarrowtail V^{I}$. Since the Yoneda embedding is full, there is an arrow $\phi\colon G\to A$ such that $\yo(\phi)=f$. Because $G$ is extremal projective, there exists an arrow $\psi\colon G\to \coprod_{I}{G}$ such that $\delta\circ \psi = \phi$. The implicit total operation $\yo(\psi)\colon V^{I}\to V$ satisfies $\yo(\psi)\circ \yo(\delta) = f$, i.e.\ it extends $f$, as was to be proved.

Conversely, suppose that every representable implicit partial operation $V^{I} \rightsquigarrow V$ with domain $\yo(\delta)\colon \yo(A)\rightarrowtail V^{I}$ can be extended to a total one. That is, for every natural transformation $f\colon \yo(A)\to V$ there exists $g\colon V^{I}\to V$ such that $g\circ \yo(\delta) = f$. Therefore, in $\cat$, for every $\phi\colon G\to A$ there exists $\psi\colon G\to \coprod_{I}G$ such that $\delta\circ \psi = \phi$. 

Let $J$ be the set $\cat(G,A)$, and let $\epsilon\colon \coprod_{J}G \to A$ be the canonical arrow. Since $G$ is an extremal generator, $\epsilon$ is an extremal epimorphism. For each $\phi\colon G \to A$, choose $\psi_{\phi}\colon G\to \coprod_{I}G$ such that $\delta\circ \psi_{\phi} = \phi$. The cocone $\{\psi_{\phi}\mid \phi\in J\}$ induces a unique mediating morphism $\iota\colon \coprod_{J}G \to \coprod_{I}G$, which satisfies $\delta\circ \iota = \epsilon$ by construction. As~$\epsilon$ is an extremal epimorphism, we conclude that $\delta$ is an extremal epimorphism. 
\end{proof}
%
It follows at once from the latter claim that \ref{i:representable-admits-extension} $\IMP$ \ref{i:ES-at-copowers-of-G}.
%
For the converse implication, let $f\colon V^{I}\rightsquigarrow V$ be an implicit partial operation whose domain $d\colon D\rightarrowtail V^{I}$ is represented by an object $A$ of $\cat$. 
%
Thus, the span in the copresheaf category $\CP{\cat}$ displayed on the left-hand side below corresponds to a cospan in $\cat$ as on the right-hand side below.
\begin{equation*}
\begin{tikzcd}
& D \arrow[rightarrowtail]{dl}[swap]{d} \arrow{dr}{f} & \\
V^{I} & & V
\end{tikzcd}
\ \ \ \ \ \ 
\begin{tikzcd}
& A  & \\
\coprod_{I}{G} \arrow{ur}{\delta} & & G \arrow{ul}[swap]{\phi}
\end{tikzcd}
\end{equation*}
%
The Yoneda embedding $\cat^{\op}\to \CP{\cat}$ reflects monomorphisms because it is faithful, so $\delta$ is an epimorphism in $\cat$. Item~\ref{i:ES-at-copowers-of-G}, combined with Claim~\ref{claim:ext-epi-iff-total}, implies that $f\colon V^{I}\rightsquigarrow V$ can be extended to an implicit total operation, as was to be~proved.
\end{proof}

\end{document}
